\clearpage
\section{Translation}
\label{sec:translation}

The calculus presented in \cref{sec:expr-proc,sec:types,sec:algorithmic-typing}
is based on  a linear and ordered type
system. In this section we first define an run-time calculus that
is closer to an implementation of splitting and borrowing in a real-world
porgramming language. On its own, the run-time calculus from this section
provides fewer guarantees. We show that translating a well typed expression
into the run-time calculus preserves any guarantees made by the
earlier type system for the source language.

\begin{figure}
  \setlength{\abovedisplayskip}{0pt}%
  \setlength{\belowdisplayskip}{0pt}%
  \begin{align*}
    \TProc \grmeq
        & \TPExp\E
          \grmor \TPPar\TProc\TProc
          \grmor \TPNew[xx] \TProc
          \grmor \TPSync \Flag \TProc
      \tag{Configurations}
    \\
    \TCC
      \grmeq& \TCCEmpty
      \grmor \TPPar \TCC \TProc
      \grmor \TPNew[xx] \TCC
      \grmor \TPSync \Flag \TCC
      \tag{Configuration contexts}
    \\
    \Flag \grmeq & \FlagInit \grmor \FlagDropped \grmor \FlagAcquired
      \tag{Synchronization flags}
  \end{align*}
  \caption{Untyped run-time syntax}
  \label{fig:unr-syntax}
\end{figure}

\begin{figure}
  \begin{mathpar}
    %\RuleTransCongRefl \and
    \RuleTransCongParComm \and
    \RuleTransCongParAssoc \and
    \RuleTransCongParUnit \and
    \RuleTransCongResSwap \and
    \RuleTransCongResComm \and
    \RuleTransCongFlagComm \and
    \RuleTransCongResFlagComm \and
    \RuleTransCongResExtend \and
    \RuleTransCongFlagExtend
    %\RuleTransCongTrans
  \end{mathpar}
  \caption{Congruence rules of translated processes}
  \label{fig:translated-congruence}
\end{figure}

The syntax of constants, values, and expressions remains unchanged.
\Cref{fig:unr-syntax} gives a new definition of process configurations
$\TProc$. Besides the embedding of expressions and the parallel
composition of processes,  $\TPNew[xy] \TProc$ introduces a channel
with endpoints $x$ and $y$ scoped over $\TProc$
and $\TPSync \Flag \TProc$ introduces a synchronization variable $z$
set to $\Flag$ scoped over $\TProc$. Channel introduction with dual
binders is standard in session-type calculi
\cite{gay.vasconcelos:session-types}; synchronization variables are
new. We call channels introduced by the run-time calculus
\emph{primitive channels} to distinguish them from translated source
channels.

Processes are subject to the usual congruence relation shown in
\cref{fig:translated-congruence}. The additional congruences deal with
flag introduction and they are analogous to the congruence rules for
channels.

% The type syntax does not differentiate between different kinds of function
% arrows. Product types are unordered. $\TBUnit$ and sum types are carried over
% unmodified. Channels are untyped and there are now two kinds: data channels
% $\TBChanD$ and synchronization channels $\TBChanF$. Communication channels
% $\TBChan$ are defined using type abbreviations:
% %
% \begin{align*}
%   \TBChanF* &= \TTSum\TBChanF\TBUnit\\
%   \TBChan &= \TTProd{\TBChanF*}{\TTProd{\TBChanD}{\TBChanF*}}
% \end{align*}

A translated expression/process represents a communication channel as
a triple $\targettermcolor{c} = 
\TEChan[\TEVar]\TEVary[\TEVarz]$ made up of an optional synchronization variable
$\TEVar$, a primitive channel $\TEVary$ and a second optional synchronization
variable $\TEVarz$. If $\TEVar$ is a synchronization, an $\EApp \CAcq c$ will have to
be executed first before any communication can happen. Payloads are transferred
via the data channel $\TEVary$. If $\TEVarz$ is an actual
synchronization, then $\targettermcolor{c}$ will have to be dropped, not terminated.

The translation eliminates bind groups $\BindGroup$. The
$\TBindNu$-binder follows the standard syntax more closely. 
The new binder $\TBindSync$ binds asynchronous channels used to synchronize
between two parts of a remote split. Synchronization must use asynchronous
communication so that a $\CDrop$ does not block until the corresponding
$\CAcq$.

\begin{figure}
  \begin{mathpar}
    \RuleUProcRedExp      \and
    \RuleUProcRedDiscard  \and
    \RuleUProcRedFork     \and
    \RuleUProcRedNew      \and
    \RuleUProcRedLSplit   \and
    \RuleUProcRedRSplit   \and
    \RuleUProcRedDrop     \and
    \RuleUProcRedAcquire  \and
    \RuleUProcRedCleanup  \and
    \RuleUProcRedClose    \and
    \RuleUProcRedCom      \and
    \RuleUProcRedChoice   \and
    \RuleUProcRedStruct   \and
    \RuleUProcRedPar      \and
    \RuleUProcRedBind     \and
    \RuleUProcRedSync
  \end{mathpar}
  \caption{Reduction of translated processes ($P' \TPReduce Q'$)}
  \label{fig:unr-process-reduction}
\end{figure}

The reduction behaviour of translated expressions is the same as their
untranslated counterpart (\cref{fig:expression-reduction}).
\Cref{fig:unr-process-reduction} defines the reduction relation for translated
processes. It uses $\ChanPat\E$ as an abbreviation for channels
$\TEChan [\E[1]] \E [\E[2]]$.

\begin{figure}
  \begin{mathpar}
    \RuleTransProcExp   \and
    \RuleTransProcPar   \and
    \RuleTransProcBind  \\
    \RuleTransBindPar   \and
    \RuleTransBindChannels \and
    \RuleTransBindSeq   \and
    \RuleTransBindVar   \and
    \RuleTransBindNil
  \end{mathpar}
  \begin{align*}
    \BindFlag{\BPar{\BindGroup[1]}{\BindGroup[2]}} &= \BindFlag{\BindGroup[2]} &
    \BindFlag{\Bind} &= \begin{cases}
      \FlagDropped & \text{if $\Fv(\Bind) = \emptyset$} \\
      \FlagInit & \text{otherwise}
    \end{cases}
  \end{align*}
  % \begin{mathpar}
  %   \inferrule{}{
  %     \Bind \BCat \BEmp = \Bind
  %   }

  %   \inferrule{
  %     \Bind[2] \ne \BEmp
  %   }{
  %     \Bind[1] \BCat \Bind[2] = \BSeq{\Bind[1]}{\Bind[2]}
  %   }

  %   \inferrule{}{
  %     \BEmp \BNorm \BEmp
  %   }

  %   \inferrule{}{
  %     \EVar \BNorm \EVar
  %   }

  %   \inferrule{
  %     \Bind[1] \BNorm \Bind*[1] \\
  %     \Bind[2] \BNorm \Bind*[2]
  %   }{
  %     \BSeq{\Bind[1]}{\Bind[2]} \BNorm \Bind*
  %   }

  %   \inferrule{
  %     \Bind \BNorm \Bind*
  %   }{
  %     \BSeq\EVar\Bind \BNorm \EVar \BCat \Bind*
  %   }
  % \end{mathpar}
  \caption{Translation of
    processes ($\Proc \TransProc \TProc*$)
    and bind groups ($\BindGroup \TransBind \TCC [\Sub]$) with $\BindFlag{}$
    mapping bind groups $\BindGroup$ to synchronization flags $\Flag$.}
  \label{fig:translation}
\end{figure}

\begin{lemma}[Bisimulation]
  \label{lemma:translation-bisimulation}
%
  For any context $\Ctx$, process $\Proc$\!, and substitution $\Sub$ if\/
    $\ProcHasType \Ctx \Proc$
  and
    $\Proc \TransProc \TProc*$\!,
  then for all
    $\ProcQ$
  where
    $\ProcQ \TransProc \TProcQ*$,
  it holds
    $\Proc \PReduce \ProcQ$
  iff
    $\TProc* \TPReduce \TProcQ*$.

  {\upshape \js{$\Sub$ needs to be restricted: should only substitute channel
  variables with values. Do we need to impose requirements on the inserted
  values?}}
\end{lemma}

%%% Local Variables:
%%% mode: latex
%%% TeX-master: "../popl2027.tex"
%%% End:
